% a01_H0rel part 5b: the lower bound
\subsection{The lower bound}\label{app:a01:lower}

Following \cite[Sect.~3]{Ba1} and \cite[Appendix]{K1}: $p_i\defeq p(\ell_i)$, $M_i\defeq p_i\lambda(\ell_i)^{-1/4}$, $M_i'\defeq p_{i-1}\lambda(\ell_{i-1})^{-1/4}=L^{1/4}(p_{i-1}/p_i)M_i$ the block-sup threshold of \cite[(3.2)]{K1}, $p_1(\ell)\defeq b_1(1+\ln\ell^{-1})^{p_1}$ the fluctuation cut-off with $p-1\ge p_1$.

\begin{definition}[the class carried through the induction]\label{def:a01:class}
In unit-$i$ units, with $\phi^{(i)}\defeq a_i\Mop(\Xl,A^{(i)})^{-1}Q(A^{(i)})^*\phi_i$, the class $(\mathrm{A.1})^{(i)}_C$ is defined by
\[
\sup_{x\in\Xl}|\phi^{(i)}(x)|\le C_1M_i,\qquad\sup_y|\phi_i(y)|\le C_1M_i,
\]
\[
(\mathrm G)_i:\ \sup_{b}|D_{A^{(i)}}\phi^{(i)}(b)|\le C_1'p_i\ \text{ over all bonds of the cube, face bonds included},\qquad(\mathrm E)_i:\ \sup_y|\phi_i(y)-Q(A^{(i)})\phi^{(i)}(y)|\le C_1''p_i,
\]
together with the photon clause $|\partial A^{(i)}|_\infty\le C'_\Gamma p_i$. On a face bond $(\mathrm G)_i$ reads $\vth^{-1}|\phi^{(i)}(x)|\le C_1'p_i$ at the adjacent live site: a fine-level linear vanishing with slope $O(p_i)$.
\end{definition}

The first two clauses are King's (A.1) with its Laplacian clause dropped; $(\mathrm G)$ and $(\mathrm E)$ are the two consequences of that clause which the lower bound consumes (\cref{lem:a01:PC}). On the cube the Laplacian clause itself holds as a consequence ($(-\Delta^D_A+V_i)\phi^{(i)}=a_iQ^*(\phi_i-Q\phi^{(i)})$ gives $|\Delta^D_A\phi^{(i)}|\le V_iC_1M_i+a_i\Wt C_1''p_i$), but it is not needed.

\begin{lemma}[compatibility, coefficient one]\label{lem:a01:GC}
Let $\psi\in(\mathrm{A.1})^{(i+1)}_C$, $|\phi_i'|_\infty\le p_1(\ell_i)$ and $\phi_i=aL^{-2}C^{(i)}Q_c^*\psi+\phi_i'$ at the gauge field of Ba{\l}aban's step ($A^{(i)}=A^{(i+1)}+A'^{(i)}$). Then $\phi_i\in(\mathrm{A.1})^{(i)}_C$ provided $L^{3/4}\ge3$, $C_1\ge3\max(C_A,1)\max(1,b_1/b_0)$, $C_1'\ge3C_D\max(1,b_1/b_0)$, $C_1''\ge\max\{2(1+C_{Q4})\max(1,b_1/b_0),C_1'\}$, and the three smallness conditions on $\eps_0^{\mathrm{Ba}}$ for the gauge-field remainders (below), where $C_A$, $C_D$ are the constants of \cref{lem:a01:A4}(b) and \cref{lem:a01:A5} for the leg of block data ($C_A,C_D\le Ca\Wt$) and $C_{Q4}\le C_A$ the constant in $|Q(A)u|_\infty\le C_{Q4}|b|_\infty$.
\end{lemma}
\begin{proof}
(1) \emph{Exact composition.} By the composition law of the kernels and the minimiser identity (the minimiser of the composed functional is the minimiser of the one-step functional with the composed average, the uniform coarse step having $Q_cQ_c^*=1$), $\phi^{(i)}(\phi_i)=\Phi_\psi+u'$ with $\Phi_\psi\defeq\phi^{(i+1)}(\psi)$ and $u'\defeq a_i\Mop^{-1}Q^*\phi_i'$, $|u'|_\infty\le C_A|\phi_i'|_\infty$. (2) \emph{Units.} A field in unit-$(i+1)$ units is $L^{1/2}$ times the field in unit-$i$ units and a length is $L$ times, so a fine-bond difference quotient converts by $L^{-3/2}$ and a block average by $L^{-1/2}$; $p_{i+1}\le p_i$, $\lambda(\ell_{i+1})=L\lambda(\ell_i)$. (3) \emph{Sup clauses.} $|\phi^{(i)}|_\infty\le|\Phi_\psi|_\infty+C_A|\phi_i'|_\infty\le C_1L^{-3/4}M_i+C_Ap_1(\ell_i)\le\frac13C_1M_i+\frac13C_1M_i$, the third share being left for the gauge-field remainder. Block clause: by the blockwise push-down identity $aL^{-2}C^{(i)}Q_c^*\psi=Q\Phi_\psi+\kappa'Q_c^*(\psi-Q_{i+1}\Phi_\psi)$ with $\kappa'=aL^{-2}/(a_i+aL^{-2})\le L^{-2}$ (minimise the composed functional over $\phi_i$ block by block: on each $(i+1)$-block the minimiser is $Q\Phi+\delta$ with $\delta$ constant on the block, and the one-dimensional minimisation gives $\delta=\kappa'(\psi-Q_{i+1}\Phi)$; this uses only that the coarse step is uniform), and $|Q\Phi|\le\sup|\Phi|$ (Q5), $|Q_c^*c|=|c|$: $|\phi_i|_\infty\le(1+2\kappa')C_1L^{-3/4}M_i+p_1(\ell_i)\le C_1M_i$. (4) $(\mathrm G)_i$ at live bonds: $|D\Phi_\psi(b)|\le L^{-3/2}C_1'p_{i+1}\le\frac13C_1'p_i$ by $(\mathrm G)_{i+1}$ at the same fine bond, and $|Du'(b)|\le C_D|\phi_i'|_\infty\le\frac13C_1'p_i$ by \cref{lem:a01:A5}; the third share is the gauge-field remainder. At face bonds the same three terms, each a face-bond derivative: $D\Phi_\psi$ by $(\mathrm G)_{i+1}$ on the same face bond, $Du'$ by the face-bond clause of \cref{lem:a01:A5}, and the remainder likewise; the difference $D_{B+A'}\Phi_\psi-D_B\Phi_\psi$ vanishes on a face bond (no link). (5) $(\mathrm E)_i$: by the push-down identity and the exact composition,
\[
\phi_i-Q(A)\phi^{(i)}(\phi_i)=\kappa'Q_c^*\big(\psi-Q_{i+1}\Phi_\psi\big)+\big(\phi_i'-Q(A)u'\big),
\]
and at every block $|\cdot|\le\kappa'L^{-1/2}C_1''p_{i+1}+(1+C_{Q4})p_1(\ell_i)\le L^{-5/2}C_1''p_i+\frac12C_1''p_i\le C_1''p_i$. (6) \emph{The gauge-field remainder.} The composition holds at a common gauge field; in Ba{\l}aban's step the push-down is taken at $B\defeq A^{(i+1)}$ while $\phi^{(i)}$ is built with $B+A'$. The background of $\psi$ at $B+A'$ differs from $\Phi_\psi$ by $\delta$ solving $\Mop(B+A')\delta=f$, $f=a_i[Q(B+A')^*-Q(B)^*]\omega_B-[\Mop(B+A')-\Mop(B)]\Phi_\psi$ with $\omega_B$ the pushed-down datum, $|\omega_B|_\infty\le C_1M_i$; expanding term by term, $f$ is a bounded function with $|f|\le Ce_ip_1(\ell_i)(C_1M_i+C_1'p_i)$ on Ba{\l}aban's fluctuation class ($|A'|,|\partial A'|\le C_\gamma p_1(\ell_i)$ by \cite[(2.18)]{K1} with \cite[Thm.~3.3]{K1} for the vector operators and the cut-off \cite[(3.50)]{Ba1}; on the cube by \cref{lem:a01:PE,lem:a01:REG}), involving only live bonds and contours inside blocks; by \cref{lem:a01:A4}(c) and \cref{lem:a01:A5} (bounded source), $|\delta|_\infty+\sup_b|D_A\delta(b)|\le C_{\mathrm{rem}}e\,p_1(\ell_i)(C_1\lambda^{-1/4}\ell_i^{1/4}+C_1'\ell_i^{1/2})p_i$, which is $\le\frac13C_1'p_i$ and $\le\frac13C_1M_i$ under $C_{\mathrm{rem}}e\,C_1\lambda^{-1/4}\beta^{1/4}p_1(\beta)\le\frac16C_1'$, $C_{\mathrm{rem}}e\,\beta^{1/2}p_1(\beta)\le\frac16$ and $C_{\mathrm{rem}}e\,p_1(\beta)(\lambda^{-1/4}\beta^{1/4}+(C_1'/C_1)\beta^{1/2})\le\frac13$ at $\beta=\eps_0^{\mathrm{Ba}}$ (monotone in $\beta$ on $(0,e^{1-4p_1}]$, so valid at every $\ell_i\le\eps_0^{\mathrm{Ba}}$); the block-averaged form of the remainder enters $(\mathrm E)_i$ with the sup clause of \cref{lem:a01:A4}(c) and is absorbed in the room $\frac12-L^{-5/2}\ge\frac13$ (here $C_1''\ge C_1'$ is used). The scale-$(i+1)$ class enters with coefficient one; no constant multiplies $p_{i+1}$, which is what \cite[(3.53)--(3.54)]{Ba1} requires, and the conditions (3.54) are unchanged.
\end{proof}

\begin{corollary}[King's face test as a consequence at every scale of the lower bound]\label{cor:a01:BR}
On $(\mathrm{A.1})^{(i+1)}_C$, at every face block $y$ of scale $i$, the push-down $b_i\defeq aL^{-2}C^{(i)}Q_c^*\psi$ obeys $|b_i(y)|\le(1+\vth)L^{-3/2}C_1'p_i+L^{-5/2}C_1''p_i$, hence $|\phi_i(y)|\le$ that $+p_1(\ell_i)$ on the induction set.
\end{corollary}
\begin{proof}
$b_i=Q\Phi_\psi+\kappa'Q_c^*(\psi-Q_{i+1}\Phi_\psi)$. $\Phi_\psi$ vanishes on the killed layer; along a normal path of length $\le1+\vth$ from the killed site to any point of $B(y)$ every bond is a live or face bond, so $|\Phi_\psi|\le(1+\vth)\sup|D\Phi_\psi|\le(1+\vth)L^{-3/2}C_1'p_i$ by $(\mathrm G)_{i+1}$ including its face part; $Q$ is an average; the second term is $(\mathrm E)_{i+1}$ with $\kappa'L^{-1/2}\le L^{-5/2}$.
\end{proof}

This is a consequence, not an inserted test: no block-field test is inserted at an intermediate scale of the lower bound (such a test would not be implied by the previous scale, the push-down norm having no numerical bound below one); King's remark that there is ``no change in the proof of the lower bound'' \cite[p.~343]{K2} is correct in this sense once King's Laplacian clause is replaced by $(\mathrm G)$, $(\mathrm E)$.

\begin{lemma}[the two properties consumed by the $n+n'=1$ vertices]\label{lem:a01:PC}
On $(\mathrm{A.1})^{(i)}_C$, at every block $y$, every live $x\in B(y)$ and every live bond $b$: (P-grad) $|D_{A^{(i)}}\phi^{(i)}(b)|\le C_1'p_i$; (P-block) $|\phi^{(i)}(x)-U(A^{(i)}(\Gamma_{x,y}))\phi_i(y)|\le C_Pp_i$ with $C_P\defeq C_1''+2dC_1'$.
\end{lemma}
\begin{proof}
(P-grad) is $(\mathrm G)_i$. (P-block): $Q\phi^{(i)}(y)=\sum_{x'}\omega(x')U(\Gamma_{y,x'})\phi^{(i)}(x')$ with $\omega\defeq\vth^3w\ge0$, $\sum\omega=1$; so $\phi^{(i)}(x)-U(\Gamma_{x,y})\phi_i(y)=U(\Gamma_{x,y})[Q\phi^{(i)}(y)-\phi_i(y)]+\sum_{x'}\omega(x')[\phi^{(i)}(x)-U(\Gamma_{x,y})U(\Gamma_{y,x'})\phi^{(i)}(x')]$; the first term is $(\mathrm E)_i$, and each term of the sum is the covariant difference along the path $\Gamma_{y,x'}$ followed by $\Gamma_{x,y}$, of length $\le2d$, the exact telescoping sum of $\vth D_A\phi^{(i)}$ over its bonds (live sites only, by the choice of contours), hence $\le2dC_1'p_i$ by $(\mathrm G)_i$.
\end{proof}

\begin{lemma}[final restriction at the stopping scale]\label{lem:a01:final}
Let $\chi_k\lambda_k$ be the small-field indicator of \cref{app:a01:weights} at the stopping scale (King's \cite[(3.2)]{K1} with the gradient clause on every block bond of the cube, times the face test). Then $\supp(\chi_k\lambda_k)\subset(\mathrm{A.1})^{(k)}_C$ once $C_1\ge\max(C_A,1)L^{1/4}(1+\ln L)^p$, $C_1'\ge C_G(1+\ln L)^p$, $C_1''\ge C_E(1+\ln L)^p$, $V_kM_k'\le p_k$ and $e_kM_k\le1$ (conditions on $\eta_1$), with $C_G,C_E$ depending on $(C_8,\delta_1,\delta',C_A,L,a,d)$ only. At the final scale $K_f$ the set $\{|\phi_{K_f}(y)|\le1\ \forall y\}$ lies in $(\mathrm{A.1})^{(K_f)}_C$.
\end{lemma}
\begin{proof}
Facts on the final set: (f1) $|\phi_k(y')|\le C_8p_k$ on $\Lone^{(k)}$; (f2) $|D_{A_k}\phi_k(\beta)|\le p_{k-1}$ on every block bond; (f3) $|\phi_k|\le M_k'$; hence (f5) $|\phi_k(y)|\le C_9(1+\dist(y,\Lone^{(k)}))p_{k-1}$, $C_9\defeq\max(C_8,1)$ (the cube is a box: every block is joined to the face layer by a cube-internal block path of that length), and (f6) $|\phi_k(y)-U(A_k(\Gamma^{\mathsf T}_{y,y_0}))\phi_k(y_0)|\le|y-y_0|_1p_{k-1}$ along a tree $\mathsf T$ of block paths. Sup clauses: (f3) and \cref{lem:a01:A4}(b), with $M_k'\le L^{1/4}(1+\ln L)^pM_k$. $(\mathrm G)_k$ at a bond $b$ (live or face) in the block $y_0$: write $\phi_k=\phi_k(y_0)\Theta^{\mathrm{blk}}+(\phi_k-\phi_k(y_0)\Theta^{\mathrm{blk}})$ with $\Theta^{\mathrm{blk}}$ the unit tree datum transported by $A_k$. Constant part: \cref{lem:a01:FL}(a) with $|b-y_0|\le1$, the tree datum transported with the block field $A_k$ differing from $\Theta_A$ by $\le Ce_kp_k|y-y_0||c|$ (the photon block-versus-background step \cite[(A.2)--(A.6)]{K1}, on the cube by \cref{lem:a01:PE,lem:a01:REG}): $\le|\phi_k(y_0)|[CV_k+Ce^{-\delta_1\dist(b,\DD)}+C\theta_k+Ce_kp_k]\le C[V_kM_k'+C_9\sup_t(1+t)e^{-\delta_1(t-1)}p_{k-1}+(\theta_k+e_kp_k)M_k']\le C'p_k$ by (f3), (f5), $\dist(b,\DD)\ge\dist(y_0,\Lone^{(k)})-1$, $V_kM_k'\le p_k$, $\theta_k\le Ce_kp_k$ and $e_kp_kM_k'\le C_1p_k$. Varying part: $\sum_y|D_{A,b}l(b,y)||\phi_k(y)-\phi_k(y_0)\Theta^{\mathrm{blk}}(y)|\le\sum_yCe^{-\delta'|b-y|}|y-y_0|_1p_{k-1}\le Cp_{k-1}$ by \cref{lem:a01:A5} in kernel form and (f6). $(\mathrm E)_k$: by \cref{lem:a01:EB}(a), $|\phi_k(y)-Q\phi^{(k)}(y)|=a_k^{-1}|\Delta^{(k)}\phi_k(y)|$; with the tree rooted at $y$, $\Delta^{(k)}\phi_k(y)=\sum_{y'}\Delta^{(k)}(y,y')(\phi_k(y')-\phi_k(y)\Theta^{\mathrm{blk}}(y'))+\phi_k(y)(\Delta^{(k)}\Theta^{\mathrm{blk}})(y)$; the kernel decays, $|\Delta^{(k)}(y,y')|\le Ce^{-\delta'|y-y'|}$ (\cref{lem:a01:A23}(b)), so the first sum is $\le Cp_{k-1}$ by (f6), and the second is $a_k|\phi_k(y)||\Theta^{\mathrm{blk}}(y)-Qu(y)|\le a_k|\phi_k(y)|[CV_k+Ce^{-\delta_1\dist(B(y),\DD)}+C(\theta_k+e_kp_k)]\le Cp_k$ by \cref{lem:a01:FL}(b), (f3), (f5). At the final scale the data are bounded by $1$ and $(\mathrm G)_{K_f}$, $(\mathrm E)_{K_f}$ follow from \cref{lem:a01:A5,lem:a01:A4}(b) with $|b|_\infty\le1$, since $p_{K_f}\ge b_0\ge1$.
\end{proof}

The face test is what places the face blocks in the class: without it, constant data of size $M_k'$ violate $(\mathrm G)_k$ near the face by the factor $M_k'/p_k$ (\cref{lem:a01:FL}(a) is sharp at the face).

\paragraph{The steps of the lower bound.} We follow \cite[Sect.~3]{Ba1} with King's modifications \cite[Appendix, p.~676]{K1}; $E_1$ is any field-independent constant (\cref{lem:a01:cnum}).
\begin{description}[leftmargin=1.5em,style=nextline]
\item[{B1 (restriction to small fields; \cite[(3.6)]{Ba1})}] The integrand of $\Zk(s)$ is positive, so inserting any indicator at any step gives a lower bound. We insert, at each step $i<k$, the indicator of the class $(\mathrm{A.1})^{(i)}_C$ (a class on the backgrounds, the block sup and the block mismatch, not King's six block-field tests), and at the stopping scale $k$ the indicator $\chi_k\lambda_k$, a subset of $(\mathrm{A.1})^{(k)}_C$ by \cref{lem:a01:final}. Nothing is paid in a lower bound.
\item[{B2 (inductive hypothesis; \cite[(3.26)--(3.32)]{Ba1})}] After $i$ steps, $\Zk\ge e^{-a_i}\int\chi_{(\mathrm{A.1})^{(i)}_C}e^{-S^{(i)}_k}$ with $a_i\defeq\sum_{i'<i}C\ell_{i'}^{\kappa'}|X^{(i')}|$ and $S^{(i)}_k$ King's effective action \cite[(3.14)]{K1} built on the cube; the identities for the normalisations and the representation of the fields after $i$ steps are \cite[(2.13)--(2.17)]{K1}, exact for every weight and boundary condition (the block-spin RG is a Gaussian identity for any positive covariance; the composition law uses only that the coarse steps are uniform).
\item[{B3 (the fluctuation step; \cite[(3.37)--(3.49), Props.~3.1--3.2]{Ba1})}] The Gaussian fluctuation integral is exact; the interaction is expanded to order $\bar n=13$ \cite[Appendix]{K1} with Cauchy estimates in the complexified couplings; the $R$-terms and the coefficient bounds are \cref{prop:a01:TD}(c) under the hypotheses supplied by \cref{def:a01:class,lem:a01:PC}; the one-loop determinants are bounded below by $\ln(1+x)\le x-x^2/2+\dots$ ($n$ odd) for the bounded operators $C^{(i)}$ of \cref{lem:a01:A23}(c).
\item[{B4 (cut-off removal and the consistency conditions; \cite[(3.50)--(3.55)]{Ba1})}] The cumulant bound $\int\chi e^Vd\mu\ge\exp[\sum_{n\le\bar n}E^T_n(V)/n!-R]$ with cut-off-free cumulants and $R=O(1)\ell_i^{\kappa'}|X^{(i)}_1|$, $\kappa'>d$, for a Gaussian measure with exponentially decaying precision (\cref{lem:a01:A1,lem:a01:A23}) and a real polynomial interaction with Ba{\l}aban's two-component cut-off: this is the lemma behind \cite[Prop.~3.2]{Ba1}, with its interaction input \cref{prop:a01:TD}(c); the compatibility inequality $\chi_{(\mathrm{A.1})^{(i+1)}_C}(\psi)\chi(|\phi_i'|\le p_1(\ell_i))\le\chi_{(\mathrm{A.1})^{(i)}_C}(\phi_i)$ is \cref{lem:a01:GC}, with coefficient one on the scale-$(i+1)$ class, the rooms $L^{3/4}$ (sup clauses), $L^{3/2}$ ($(\mathrm G)$), $L^{5/2}$ ($(\mathrm E)$), and the conditions \cite[(3.54)]{Ba1} unchanged.
\item[{B5 (King's Appendix: \cite[(3.2)]{K1} in place of \cite[(3.27)--(3.29)]{Ba1})}] King's Laplacian clause is replaced by $(\mathrm G)_i$, $(\mathrm E)_i$, which compose with coefficient one (\cref{lem:a01:GC}) and supply through \cref{lem:a01:PC} exactly the two properties of the external scalar field that \cite[p.~421]{Ba3} draws from \cite[Lemma~2.4]{Ba2} for his class; what \cite[Prop.~1]{Ba3} uses of the external legs otherwise is their norm, \cref{lem:a01:legs}(a). The final set $\chi_k\lambda_k$ is in the class by \cref{lem:a01:final}. The photon background's clauses \cite[(A.2)--(A.6)]{K1} are unchanged on the cube (\cref{app:a01:photon}).
\item[{B6 (the final step at $\ell_f=\ell_{K_f}$; \cite[(3.66)--(3.69)]{Ba1})}] Expand in the small $A^{(K_f)}$ (\cref{prop:a01:TD}(c)); $|V^{(K_f)}|\le O(1)\ell_f^{\beta_0}|X^{(K_f)}|$; the finite-dimensional integral over $\{|\phi_{K_f}(y)|\le1,|A_{K_f}|\le1\}\subset(\mathrm{A.1})^{(K_f)}_C$ (\cref{lem:a01:final}) of $\exp[-\frac12\ip{A_{K_f}}{\Delta^{(K_f)}A_{K_f}}-\frac12\ip{\phi_{K_f}}{\Delta^{(K_f)}\phi_{K_f}}]$ is $\ge\exp(-O(1)|X^{(K_f)}|)$, $|X^{(K_f)}|=\ell_f^{-d}|\CJ|$, using only the upper bounds $\Delta^{(K_f)}\le a_{K_f}\ell_f^{-2}$ on the integration region; and the normalisation identity \cite[(3.69)]{Ba1} gives $Z_{K_f}Z_{K_f}(0)e^{-E_0}\ge\exp[-\frac12a_{K_f}\ell_f^{d-2}(\Tr P_{K_f}G+\Tr P_{K_f}G^\gamma)]$ with $\Tr P_{K_f}G\le W_2\Tr(\Pi G)\le W_2N|X^{(K_f)}|\norm G$ (Q3). Hence $\Zk(0)\ge e^{-E_-|\CJ|}$ in the normalisation in which $E_1$ is the domain-adapted constant of \cite[(3.35)]{K1} on $\CJ$, with $E_-=O(\ell_f^{-d})$ depending on $\eps_0^{\mathrm{Ba}}$ and King's data only.
\item[B7 (the source)] $F(u_x)$ is King's plaquette source \cite[(2.7)]{K1}, whose vertices \cite[(3.20)]{K1} sit at $\ell^\infty$-distance $>1$ from $\R^3\setminus\CJ$ (zone I), so the face enters them only through propagators (\cref{prop:a01:TD}).
\end{description}

\begin{proposition}[lower line]\label{prop:a01:lower}
For $\eta\le\eta_1$ (so that \cite[(3.54)]{Ba1} holds at all scales $\le k$ and $j\le k-J_*$) and $j\ge J_*$: $\Zk(s)\ge e^{-a}\Ik j(s)$ with $a=\sum_{i<j}C\ell_i^{\kappa'}|X^{(i)}|=C\eta^\sigma|\CJ|$, $\sigma=\kappa'-d>0$; and $\Zk(0)\ge e^{-E_-|\CJ|}$ in the normalisation of B6.
\end{proposition}
\begin{proof}
B1--B7. The sum $\sum_{i<j}\ell_i^{\kappa'}|X^{(i)}|=|\CJ|\sum_{i<j}\ell_i^{\kappa'-d}$ is geometric and dominated by $i=j-1$, giving $a=C\eta^\sigma|\CJ|$ with $C$ independent of $j$ and $k$ (\cref{prop:a01:TD}, \cref{lem:a01:GC}, whose constants are free of $k$, $j$, $\vth$); the final restriction is a positivity step inside the class (\cref{lem:a01:final}), so the threshold $\sigma_k$ never enters a constant.
\end{proof}
